\chapter{Exact closure and welfare decompositions}
\label{ch:decomposition}

\section{Common-baseline comparisons}

An extended welfare elasticity can differ from traffic for several reasons.
To separate them, every comparison must use the same observed baseline and the
same policy forcing. Recalibrating each alternative closure would mix the
mechanism of interest with a change in starting values.

The package constructs six objects:
\begin{longtable}{p{0.10\textwidth}p{0.22\textwidth}p{0.25\textwidth}p{0.28\textwidth}}
\toprule
Closure & Routes & Modes & Congestion \\
\midrule
\endhead
\code{H} & Not solved & Not solved & Traffic-only benchmark \\
\code{NC} & Flexible & Flexible & None \\
\code{NT} & Flexible & Flexible & Edge congestion \\
\code{F} & Flexible & Flexible & All declared channels \\
\code{FM} & Flexible & Fixed at baseline shares & All declared channels \\
\code{FR} & Baseline OD-edge incidence fixed & Flexible & All declared channels \\
\bottomrule
\end{longtable}

The labels describe local derivatives, not finite counterfactual models. In
particular, \code{FR} holds the differentiated route incidence fixed while
retaining modal response and the remaining equilibrium system.

\section{Normalized multipliers}

For a realized-friction shock on edge \(e\), define
\[
 m_{e,S}=\frac{\Ereal_{e,S}}{\rho\Xi_e}.
\]
The closures obey
\[
 m_F=m_{NC}-d_{edge}-d_{terminal}
    =m_{FM}+d_{mode}
    =m_{FR}+d_{route}.
\]
The terms are signed. A positive road wedge in multiplier units need not imply
a positive contribution to the welfare elasticity when \(\rho<0\).

The first equality gives an additive congestion ladder. The other two compare
the full model separately with fixed modes and fixed routes. Mode and route
wedges are not added to the congestion ladder because each is an alternative
closure comparison.

\section{Inverse-gap identity}

For two Jacobians \(J_A\) and \(J_B\), welfare row \(q\), and forcing \(b\),
\[
 q^\top J_A^{-1}b-q^\top J_B^{-1}b
 =q^\top J_A^{-1}(J_B-J_A)J_B^{-1}b.
\]
The package checks this identity for the edge, terminal, mode, and route
comparisons. It is a direct algebraic test of the signs, orientation, and
closure construction.

\section{Analytical Jacobian blocks}

Each closure Jacobian is represented as
\[
 J_S=D_S+u_Sv_S^\top
     +\Delta_{edge,S}
     +\Delta_{terminal,S}.
\]
The sparse block \(D_S\) contains local allocation and route shares. The
rank-one term imposes aggregate labor feedback and normalization. The last
two terms are analytically constructed congestion updates.

For road and terminal rungs the inverse difference is evaluated with a
structured parallel-sum formula. Their allocation and equilibrium-correction
components are zero under the common-baseline construction; the scarcity
component carries the wedge.

For mode and route comparisons, the package evaluates the mixed update
\[
 \widetilde K_{A:F}=J_F+S_{A:F}+C_{A:F},
 \qquad
 J_A=\widetilde K_{A:F}+R_{A:F}Q_{A:F}^\top.
\]
The welfare gap is divided into:
\[
 d_r=d_r^{alloc}+d_r^{scar}+d_r^{eq},
 \qquad r\in\{mode,route\}.
\]
The allocation term changes the sparse core, the scarcity term changes
congestion exposure, and the equilibrium term is the Woodbury correction
associated with aggregate feedback. In the current common-baseline model, the
sparse and aggregate-feedback terms are common across transport closures, so
some channels evaluate to exact structural zeros. The code evaluates the
general formulas and checks the reconstruction rather than writing zeros by
convention.

\section{The Hulten gap}

The realized-friction comparison is
\[
 \Xi_e-\Ereal_{e,F}
 =\Xi_e(1-\rho)+\rho\Xi_e(1-m_{e,F}).
\]
The first term is the common externality scaling. The second is equilibrium
propagation relative to the traffic-only benchmark.

Using the congestion ladder and primitive pass-through gives
\begin{align*}
 \Xi_e-\Eprimitive_{e,F}
 ={}&\Xi_e(1-\rho)
 +\rho\Xi_e(1-m_{e,NC})\\
 &+\rho\Xi_e d_{edge,e}
 +\rho\Xi_e d_{terminal,e}
 +\rho\Xi_e m_{e,F}(1-\chi_e).
\end{align*}
The output reports each term in levels. Signed percentages can be misleading
when large positive and negative components nearly cancel.

\section{Physical-link aggregation}

The model is directed. To obtain a two-direction physical-link experiment, the
software first sums:
\begin{itemize}
  \item primitive and realized elasticities;
  \item Hulten traffic;
  \item each additive Hulten-gap component;
  \item each analytical channel contribution.
\end{itemize}
Only after summation does it calculate a physical-link normalized multiplier.
This order avoids averaging ratios with different traffic denominators.

\warningbox{A physical-link result is defined only when exactly two declared
policy arcs have opposite endpoints and a common physical-link identifier.
Transit services and one-way links that lack such a pair remain directed
policies.}

\section{Reading a decomposition}

A useful interpretation sequence is:
\begin{enumerate}
  \item compare the traditional statistic with the full primitive result;
  \item separate common externality scaling from link-specific propagation;
  \item inspect edge and terminal congestion;
  \item compare primitive and realized effects to measure pass-through;
  \item use mode and route closures to diagnose ranking changes;
  \item check identity residuals before discussing economics.
\end{enumerate}
The decomposition explains a result conditional on the model. It does not
identify which parameter or closure is empirically correct.
